% =====================================================================
%  minggu-07.tex  --  Slide Pertemuan Minggu ke-7
%  Pengantar Kriptografi (MPFA014) -- Prodi Matematika FMIPA UBB
%
%  Topik silabus : Kriptografi modern dan stream cipher
%  Sumber materi : buku_ajar/chapters/05 (Bab 5, \S1--\S3 dan contoh)
%  Sub-CPMK      : 1.2
%
%  Kompilasi dari dalam folder ppt/:
%      pdflatex -interaction=nonstopmode minggu-07.tex   (dua kali)
% =====================================================================

\documentclass[aspectratio=169,11pt]{beamer}
% =====================================================================
%  beamer-preamble.tex  --  prasarana dan gaya slide kuliah
%
%  Mata kuliah : Pengantar Kriptografi (MPFA014)
%  Program     : Program Studi Matematika, FMIPA, Universitas Bale Bandung
%
%  Berkas ini TIDAK memuat \documentclass. Setiap berkas minggu memuatnya
%  sendiri lalu menarik berkas ini:
%
%      \documentclass[aspectratio=169,11pt]{beamer}
%      % =====================================================================
%  beamer-preamble.tex  --  prasarana dan gaya slide kuliah
%
%  Mata kuliah : Pengantar Kriptografi (MPFA014)
%  Program     : Program Studi Matematika, FMIPA, Universitas Bale Bandung
%
%  Berkas ini TIDAK memuat \documentclass. Setiap berkas minggu memuatnya
%  sendiri lalu menarik berkas ini:
%
%      \documentclass[aspectratio=169,11pt]{beamer}
%      % =====================================================================
%  beamer-preamble.tex  --  prasarana dan gaya slide kuliah
%
%  Mata kuliah : Pengantar Kriptografi (MPFA014)
%  Program     : Program Studi Matematika, FMIPA, Universitas Bale Bandung
%
%  Berkas ini TIDAK memuat \documentclass. Setiap berkas minggu memuatnya
%  sendiri lalu menarik berkas ini:
%
%      \documentclass[aspectratio=169,11pt]{beamer}
%      \input{beamer-preamble}
%      \title{...} \subtitle{...} \setWeekInfo{Minggu n}{Topik}
%
%  Berkas minggu hanya berisi \setWeekInfo dan lingkungan frame.
%  Kompilasi dijalankan DARI DALAM folder ppt/ (lihat compile-ppt.bat),
%  karena gambar buku dirujuk dengan jalur relatif ../buku_ajar/.
% =====================================================================

% ---------------------------------------------------------------- bahasa
\usepackage[utf8]{inputenc}
\usepackage[T1]{fontenc}
\usepackage[indonesian]{babel}
\usepackage{lmodern}
\usepackage{microtype}

% ------------------------------------------------------------- matematika
\usepackage{amsmath,amssymb,mathtools}

% ----------------------------------------------------------------- tabel
\usepackage{booktabs}
\usepackage{array}
\usepackage{tabularx}

% --------------------------------------------------------------- grafik
\usepackage{tikz}
\usetikzlibrary{arrows.meta,positioning,shapes.geometric,calc,fit,backgrounds}

% ---------------------------------------------------------- kode program
\usepackage{listings}
\lstdefinestyle{slidecode}{
    basicstyle=\ttfamily\scriptsize,
    keywordstyle=\color{ubbblue}\bfseries,
    commentstyle=\color{ubbgreen}\itshape,
    stringstyle=\color{ubbred},
    numbers=left,
    numberstyle=\tiny\color{gray},
    numbersep=6pt,
    frame=single,
    rulecolor=\color{ubbblue!40},
    backgroundcolor=\color{ubbblue!3},
    breaklines=true,
    showstringspaces=false,
    tabsize=4,
    columns=fullflexible,
    keepspaces=true,
}
\lstset{style=slidecode}

% =====================================================================
%  WARNA INSTITUSI
% =====================================================================
\definecolor{ubbblue}{RGB}{12, 62, 120}      % warna utama
\definecolor{ubbteal}{RGB}{0, 118, 110}      % aksen kedua
\definecolor{ubbgreen}{RGB}{22, 122, 74}     % untuk blok contoh
\definecolor{ubbred}{RGB}{168, 42, 42}       % untuk blok peringatan
\definecolor{ubbgray}{RGB}{70, 76, 84}       % teks sekunder

% =====================================================================
%  IDENTITAS MATA KULIAH
% =====================================================================
\newcommand{\CourseName}{Pengantar Kriptografi}
\newcommand{\CourseCode}{MPFA014}
\newcommand{\CourseSemester}{Semester Ganjil}
\newcommand{\CourseProgram}{Program Studi Matematika}
\newcommand{\CourseFaculty}{Fakultas Matematika dan Ilmu Pengetahuan Alam}
\newcommand{\CourseInstitution}{Universitas Bale Bandung}
\newcommand{\CourseLecturer}{Cecep Suwanda, S.Si., M.Kom.}

% Label minggu, disetel tiap berkas lewat \setWeekInfo
\newcommand{\WeekLabel}{}
\newcommand{\WeekTopic}{}
\newcommand{\setWeekInfo}[2]{%
    \renewcommand{\WeekLabel}{#1}%
    \renewcommand{\WeekTopic}{#2}%
}

% =====================================================================
%  TEMA: header, footer, warna blok
% =====================================================================
\usetheme{default}
\setbeamertemplate{navigation symbols}{}
\setbeamertemplate{blocks}[rounded][shadow=false]

\setbeamercolor{palette primary}{bg=ubbblue, fg=white}
\setbeamercolor{palette secondary}{bg=ubbgray, fg=white}
\setbeamercolor{palette tertiary}{bg=ubbteal, fg=white}

\setbeamercolor{structure}{fg=ubbblue}
\setbeamercolor{normal text}{fg=black!92}
\setbeamercolor{alerted text}{fg=ubbred}

\setbeamercolor{block title}{bg=ubbblue, fg=white}
\setbeamercolor{block body}{bg=ubbblue!6, fg=black!92}
\setbeamercolor{block title alerted}{bg=ubbred, fg=white}
\setbeamercolor{block body alerted}{bg=ubbred!6, fg=black!92}
\setbeamercolor{block title example}{bg=ubbgreen, fg=white}
\setbeamercolor{block body example}{bg=ubbgreen!6, fg=black!92}

\setbeamercolor{itemize item}{fg=ubbblue}
\setbeamercolor{itemize subitem}{fg=ubbteal}
\setbeamercolor{enumerate item}{fg=ubbblue}

\setbeamerfont{headline}{size=\scriptsize}
\setbeamerfont{footline}{size=\scriptsize}
\setbeamerfont{block title}{size=\normalsize}

% Badan blok rata kiri (bukan rata kanan-kiri): teks slide lebih mudah dibaca
% dan baris terakhir yang pendek tidak lagi memicu "Underfull \hbox".
\addtobeamertemplate{block body begin}{\raggedright}{}
\addtobeamertemplate{block alerted body begin}{\raggedright}{}
\addtobeamertemplate{block example body begin}{\raggedright}{}

% Header: nama mata kuliah dan kode MK
\setbeamertemplate{headline}{%
    \begin{beamercolorbox}[wd=\paperwidth,ht=2.6ex,dp=1.2ex,leftskip=1em,rightskip=1em]{palette primary}%
        \usebeamerfont{headline}\CourseName\ \textperiodcentered\ \CourseCode%
        \hfill \CourseProgram%
    \end{beamercolorbox}%
}

% Footer: minggu dan nomor halaman
\setbeamertemplate{footline}{%
    \begin{beamercolorbox}[wd=\paperwidth,ht=2.4ex,dp=1.2ex,leftskip=1em,rightskip=1em]{palette secondary}%
        \usebeamerfont{footline}\WeekLabel\ \textperiodcentered\ \WeekTopic%
        \hfill \insertframenumber\,/\,\inserttotalframenumber%
    \end{beamercolorbox}%
}

% Halaman judul: judul, subjudul, lalu blok identitas mata kuliah
\setbeamertemplate{title page}{%
    \begin{centering}
        \vskip1.0em
        {\usebeamerfont{title}\usebeamercolor[fg]{title}\inserttitle\par}
        \vskip0.3em
        {\usebeamerfont{subtitle}\usebeamercolor[fg]{subtitle}\insertsubtitle\par}
        \vskip1.2em
        \textcolor{ubbblue}{\rule{0.72\textwidth}{0.8pt}}
        \vskip1.2em
        {\small
            \textbf{\CourseName}\ (\CourseCode)\par
            \vskip0.4em
            \CourseProgram\par
            \CourseFaculty\par
            \CourseInstitution\par
            \vskip0.6em
            \CourseSemester\ \textperiodcentered\ \WeekLabel\par
            \vskip0.8em
            Dosen Pengampu: \CourseLecturer\par
        }
        \vskip1.0em
        \textcolor{ubbblue}{\rule{0.72\textwidth}{0.4pt}}
    \end{centering}
}

% =====================================================================
%  PERINTAH BANTU
% =====================================================================
\newcommand{\R}{\mathbb{R}}
\newcommand{\N}{\mathbb{N}}
\newcommand{\Z}{\mathbb{Z}}
\newcommand{\code}[1]{\texttt{#1}}

% Rujukan gambar dari buku: \bookfig[lebar]{NN}{namaberkas.png}
% Jalur relatif terhadap folder ppt/ (dikompilasi dari sana).
\newcommand{\bookfig}[3][0.80\linewidth]{%
    \includegraphics[width=#1]{../buku_ajar/chapters/#2/figures/#3}%
}

% Keterangan sumber gambar, sepadan dengan \sumber di buku
\newcommand{\figsumber}[1]{%
    \par\vspace{1pt}{\tiny\itshape Sumber: #1\par}%
}

% Penanda blok materi yang sering dipakai di slide
\newcommand{\slidejudul}[1]{{\large\bfseries\textcolor{ubbblue}{#1}}\par\vspace{2pt}}
\newcommand{\kunci}[1]{\textcolor{ubbred}{\textbf{#1}}}

%      \title{...} \subtitle{...} \setWeekInfo{Minggu n}{Topik}
%
%  Berkas minggu hanya berisi \setWeekInfo dan lingkungan frame.
%  Kompilasi dijalankan DARI DALAM folder ppt/ (lihat compile-ppt.bat),
%  karena gambar buku dirujuk dengan jalur relatif ../buku_ajar/.
% =====================================================================

% ---------------------------------------------------------------- bahasa
\usepackage[utf8]{inputenc}
\usepackage[T1]{fontenc}
\usepackage[indonesian]{babel}
\usepackage{lmodern}
\usepackage{microtype}

% ------------------------------------------------------------- matematika
\usepackage{amsmath,amssymb,mathtools}

% ----------------------------------------------------------------- tabel
\usepackage{booktabs}
\usepackage{array}
\usepackage{tabularx}

% --------------------------------------------------------------- grafik
\usepackage{tikz}
\usetikzlibrary{arrows.meta,positioning,shapes.geometric,calc,fit,backgrounds}

% ---------------------------------------------------------- kode program
\usepackage{listings}
\lstdefinestyle{slidecode}{
    basicstyle=\ttfamily\scriptsize,
    keywordstyle=\color{ubbblue}\bfseries,
    commentstyle=\color{ubbgreen}\itshape,
    stringstyle=\color{ubbred},
    numbers=left,
    numberstyle=\tiny\color{gray},
    numbersep=6pt,
    frame=single,
    rulecolor=\color{ubbblue!40},
    backgroundcolor=\color{ubbblue!3},
    breaklines=true,
    showstringspaces=false,
    tabsize=4,
    columns=fullflexible,
    keepspaces=true,
}
\lstset{style=slidecode}

% =====================================================================
%  WARNA INSTITUSI
% =====================================================================
\definecolor{ubbblue}{RGB}{12, 62, 120}      % warna utama
\definecolor{ubbteal}{RGB}{0, 118, 110}      % aksen kedua
\definecolor{ubbgreen}{RGB}{22, 122, 74}     % untuk blok contoh
\definecolor{ubbred}{RGB}{168, 42, 42}       % untuk blok peringatan
\definecolor{ubbgray}{RGB}{70, 76, 84}       % teks sekunder

% =====================================================================
%  IDENTITAS MATA KULIAH
% =====================================================================
\newcommand{\CourseName}{Pengantar Kriptografi}
\newcommand{\CourseCode}{MPFA014}
\newcommand{\CourseSemester}{Semester Ganjil}
\newcommand{\CourseProgram}{Program Studi Matematika}
\newcommand{\CourseFaculty}{Fakultas Matematika dan Ilmu Pengetahuan Alam}
\newcommand{\CourseInstitution}{Universitas Bale Bandung}
\newcommand{\CourseLecturer}{Cecep Suwanda, S.Si., M.Kom.}

% Label minggu, disetel tiap berkas lewat \setWeekInfo
\newcommand{\WeekLabel}{}
\newcommand{\WeekTopic}{}
\newcommand{\setWeekInfo}[2]{%
    \renewcommand{\WeekLabel}{#1}%
    \renewcommand{\WeekTopic}{#2}%
}

% =====================================================================
%  TEMA: header, footer, warna blok
% =====================================================================
\usetheme{default}
\setbeamertemplate{navigation symbols}{}
\setbeamertemplate{blocks}[rounded][shadow=false]

\setbeamercolor{palette primary}{bg=ubbblue, fg=white}
\setbeamercolor{palette secondary}{bg=ubbgray, fg=white}
\setbeamercolor{palette tertiary}{bg=ubbteal, fg=white}

\setbeamercolor{structure}{fg=ubbblue}
\setbeamercolor{normal text}{fg=black!92}
\setbeamercolor{alerted text}{fg=ubbred}

\setbeamercolor{block title}{bg=ubbblue, fg=white}
\setbeamercolor{block body}{bg=ubbblue!6, fg=black!92}
\setbeamercolor{block title alerted}{bg=ubbred, fg=white}
\setbeamercolor{block body alerted}{bg=ubbred!6, fg=black!92}
\setbeamercolor{block title example}{bg=ubbgreen, fg=white}
\setbeamercolor{block body example}{bg=ubbgreen!6, fg=black!92}

\setbeamercolor{itemize item}{fg=ubbblue}
\setbeamercolor{itemize subitem}{fg=ubbteal}
\setbeamercolor{enumerate item}{fg=ubbblue}

\setbeamerfont{headline}{size=\scriptsize}
\setbeamerfont{footline}{size=\scriptsize}
\setbeamerfont{block title}{size=\normalsize}

% Badan blok rata kiri (bukan rata kanan-kiri): teks slide lebih mudah dibaca
% dan baris terakhir yang pendek tidak lagi memicu "Underfull \hbox".
\addtobeamertemplate{block body begin}{\raggedright}{}
\addtobeamertemplate{block alerted body begin}{\raggedright}{}
\addtobeamertemplate{block example body begin}{\raggedright}{}

% Header: nama mata kuliah dan kode MK
\setbeamertemplate{headline}{%
    \begin{beamercolorbox}[wd=\paperwidth,ht=2.6ex,dp=1.2ex,leftskip=1em,rightskip=1em]{palette primary}%
        \usebeamerfont{headline}\CourseName\ \textperiodcentered\ \CourseCode%
        \hfill \CourseProgram%
    \end{beamercolorbox}%
}

% Footer: minggu dan nomor halaman
\setbeamertemplate{footline}{%
    \begin{beamercolorbox}[wd=\paperwidth,ht=2.4ex,dp=1.2ex,leftskip=1em,rightskip=1em]{palette secondary}%
        \usebeamerfont{footline}\WeekLabel\ \textperiodcentered\ \WeekTopic%
        \hfill \insertframenumber\,/\,\inserttotalframenumber%
    \end{beamercolorbox}%
}

% Halaman judul: judul, subjudul, lalu blok identitas mata kuliah
\setbeamertemplate{title page}{%
    \begin{centering}
        \vskip1.0em
        {\usebeamerfont{title}\usebeamercolor[fg]{title}\inserttitle\par}
        \vskip0.3em
        {\usebeamerfont{subtitle}\usebeamercolor[fg]{subtitle}\insertsubtitle\par}
        \vskip1.2em
        \textcolor{ubbblue}{\rule{0.72\textwidth}{0.8pt}}
        \vskip1.2em
        {\small
            \textbf{\CourseName}\ (\CourseCode)\par
            \vskip0.4em
            \CourseProgram\par
            \CourseFaculty\par
            \CourseInstitution\par
            \vskip0.6em
            \CourseSemester\ \textperiodcentered\ \WeekLabel\par
            \vskip0.8em
            Dosen Pengampu: \CourseLecturer\par
        }
        \vskip1.0em
        \textcolor{ubbblue}{\rule{0.72\textwidth}{0.4pt}}
    \end{centering}
}

% =====================================================================
%  PERINTAH BANTU
% =====================================================================
\newcommand{\R}{\mathbb{R}}
\newcommand{\N}{\mathbb{N}}
\newcommand{\Z}{\mathbb{Z}}
\newcommand{\code}[1]{\texttt{#1}}

% Rujukan gambar dari buku: \bookfig[lebar]{NN}{namaberkas.png}
% Jalur relatif terhadap folder ppt/ (dikompilasi dari sana).
\newcommand{\bookfig}[3][0.80\linewidth]{%
    \includegraphics[width=#1]{../buku_ajar/chapters/#2/figures/#3}%
}

% Keterangan sumber gambar, sepadan dengan \sumber di buku
\newcommand{\figsumber}[1]{%
    \par\vspace{1pt}{\tiny\itshape Sumber: #1\par}%
}

% Penanda blok materi yang sering dipakai di slide
\newcommand{\slidejudul}[1]{{\large\bfseries\textcolor{ubbblue}{#1}}\par\vspace{2pt}}
\newcommand{\kunci}[1]{\textcolor{ubbred}{\textbf{#1}}}

%      \title{...} \subtitle{...} \setWeekInfo{Minggu n}{Topik}
%
%  Berkas minggu hanya berisi \setWeekInfo dan lingkungan frame.
%  Kompilasi dijalankan DARI DALAM folder ppt/ (lihat compile-ppt.bat),
%  karena gambar buku dirujuk dengan jalur relatif ../buku_ajar/.
% =====================================================================

% ---------------------------------------------------------------- bahasa
\usepackage[utf8]{inputenc}
\usepackage[T1]{fontenc}
\usepackage[indonesian]{babel}
\usepackage{lmodern}
\usepackage{microtype}

% ------------------------------------------------------------- matematika
\usepackage{amsmath,amssymb,mathtools}

% ----------------------------------------------------------------- tabel
\usepackage{booktabs}
\usepackage{array}
\usepackage{tabularx}

% --------------------------------------------------------------- grafik
\usepackage{tikz}
\usetikzlibrary{arrows.meta,positioning,shapes.geometric,calc,fit,backgrounds}

% ---------------------------------------------------------- kode program
\usepackage{listings}
\lstdefinestyle{slidecode}{
    basicstyle=\ttfamily\scriptsize,
    keywordstyle=\color{ubbblue}\bfseries,
    commentstyle=\color{ubbgreen}\itshape,
    stringstyle=\color{ubbred},
    numbers=left,
    numberstyle=\tiny\color{gray},
    numbersep=6pt,
    frame=single,
    rulecolor=\color{ubbblue!40},
    backgroundcolor=\color{ubbblue!3},
    breaklines=true,
    showstringspaces=false,
    tabsize=4,
    columns=fullflexible,
    keepspaces=true,
}
\lstset{style=slidecode}

% =====================================================================
%  WARNA INSTITUSI
% =====================================================================
\definecolor{ubbblue}{RGB}{12, 62, 120}      % warna utama
\definecolor{ubbteal}{RGB}{0, 118, 110}      % aksen kedua
\definecolor{ubbgreen}{RGB}{22, 122, 74}     % untuk blok contoh
\definecolor{ubbred}{RGB}{168, 42, 42}       % untuk blok peringatan
\definecolor{ubbgray}{RGB}{70, 76, 84}       % teks sekunder

% =====================================================================
%  IDENTITAS MATA KULIAH
% =====================================================================
\newcommand{\CourseName}{Pengantar Kriptografi}
\newcommand{\CourseCode}{MPFA014}
\newcommand{\CourseSemester}{Semester Ganjil}
\newcommand{\CourseProgram}{Program Studi Matematika}
\newcommand{\CourseFaculty}{Fakultas Matematika dan Ilmu Pengetahuan Alam}
\newcommand{\CourseInstitution}{Universitas Bale Bandung}
\newcommand{\CourseLecturer}{Cecep Suwanda, S.Si., M.Kom.}

% Label minggu, disetel tiap berkas lewat \setWeekInfo
\newcommand{\WeekLabel}{}
\newcommand{\WeekTopic}{}
\newcommand{\setWeekInfo}[2]{%
    \renewcommand{\WeekLabel}{#1}%
    \renewcommand{\WeekTopic}{#2}%
}

% =====================================================================
%  TEMA: header, footer, warna blok
% =====================================================================
\usetheme{default}
\setbeamertemplate{navigation symbols}{}
\setbeamertemplate{blocks}[rounded][shadow=false]

\setbeamercolor{palette primary}{bg=ubbblue, fg=white}
\setbeamercolor{palette secondary}{bg=ubbgray, fg=white}
\setbeamercolor{palette tertiary}{bg=ubbteal, fg=white}

\setbeamercolor{structure}{fg=ubbblue}
\setbeamercolor{normal text}{fg=black!92}
\setbeamercolor{alerted text}{fg=ubbred}

\setbeamercolor{block title}{bg=ubbblue, fg=white}
\setbeamercolor{block body}{bg=ubbblue!6, fg=black!92}
\setbeamercolor{block title alerted}{bg=ubbred, fg=white}
\setbeamercolor{block body alerted}{bg=ubbred!6, fg=black!92}
\setbeamercolor{block title example}{bg=ubbgreen, fg=white}
\setbeamercolor{block body example}{bg=ubbgreen!6, fg=black!92}

\setbeamercolor{itemize item}{fg=ubbblue}
\setbeamercolor{itemize subitem}{fg=ubbteal}
\setbeamercolor{enumerate item}{fg=ubbblue}

\setbeamerfont{headline}{size=\scriptsize}
\setbeamerfont{footline}{size=\scriptsize}
\setbeamerfont{block title}{size=\normalsize}

% Badan blok rata kiri (bukan rata kanan-kiri): teks slide lebih mudah dibaca
% dan baris terakhir yang pendek tidak lagi memicu "Underfull \hbox".
\addtobeamertemplate{block body begin}{\raggedright}{}
\addtobeamertemplate{block alerted body begin}{\raggedright}{}
\addtobeamertemplate{block example body begin}{\raggedright}{}

% Header: nama mata kuliah dan kode MK
\setbeamertemplate{headline}{%
    \begin{beamercolorbox}[wd=\paperwidth,ht=2.6ex,dp=1.2ex,leftskip=1em,rightskip=1em]{palette primary}%
        \usebeamerfont{headline}\CourseName\ \textperiodcentered\ \CourseCode%
        \hfill \CourseProgram%
    \end{beamercolorbox}%
}

% Footer: minggu dan nomor halaman
\setbeamertemplate{footline}{%
    \begin{beamercolorbox}[wd=\paperwidth,ht=2.4ex,dp=1.2ex,leftskip=1em,rightskip=1em]{palette secondary}%
        \usebeamerfont{footline}\WeekLabel\ \textperiodcentered\ \WeekTopic%
        \hfill \insertframenumber\,/\,\inserttotalframenumber%
    \end{beamercolorbox}%
}

% Halaman judul: judul, subjudul, lalu blok identitas mata kuliah
\setbeamertemplate{title page}{%
    \begin{centering}
        \vskip1.0em
        {\usebeamerfont{title}\usebeamercolor[fg]{title}\inserttitle\par}
        \vskip0.3em
        {\usebeamerfont{subtitle}\usebeamercolor[fg]{subtitle}\insertsubtitle\par}
        \vskip1.2em
        \textcolor{ubbblue}{\rule{0.72\textwidth}{0.8pt}}
        \vskip1.2em
        {\small
            \textbf{\CourseName}\ (\CourseCode)\par
            \vskip0.4em
            \CourseProgram\par
            \CourseFaculty\par
            \CourseInstitution\par
            \vskip0.6em
            \CourseSemester\ \textperiodcentered\ \WeekLabel\par
            \vskip0.8em
            Dosen Pengampu: \CourseLecturer\par
        }
        \vskip1.0em
        \textcolor{ubbblue}{\rule{0.72\textwidth}{0.4pt}}
    \end{centering}
}

% =====================================================================
%  PERINTAH BANTU
% =====================================================================
\newcommand{\R}{\mathbb{R}}
\newcommand{\N}{\mathbb{N}}
\newcommand{\Z}{\mathbb{Z}}
\newcommand{\code}[1]{\texttt{#1}}

% Rujukan gambar dari buku: \bookfig[lebar]{NN}{namaberkas.png}
% Jalur relatif terhadap folder ppt/ (dikompilasi dari sana).
\newcommand{\bookfig}[3][0.80\linewidth]{%
    \includegraphics[width=#1]{../buku_ajar/chapters/#2/figures/#3}%
}

% Keterangan sumber gambar, sepadan dengan \sumber di buku
\newcommand{\figsumber}[1]{%
    \par\vspace{1pt}{\tiny\itshape Sumber: #1\par}%
}

% Penanda blok materi yang sering dipakai di slide
\newcommand{\slidejudul}[1]{{\large\bfseries\textcolor{ubbblue}{#1}}\par\vspace{2pt}}
\newcommand{\kunci}[1]{\textcolor{ubbred}{\textbf{#1}}}


\title{Kriptografi Modern dan Stream Cipher}
\subtitle{Dari bit dan XOR sampai RC4, A5/1, Trivium, Salsa20, dan ChaCha20}
\setWeekInfo{Minggu 7}{Kriptografi modern dan stream cipher}

\begin{document}

% ----------------------------------------------------------- halaman judul
\begin{frame}[plain]
    \titlepage
\end{frame}

% ------------------------------------------------------------------ peta
\begin{frame}{Peta Materi Pertemuan Ini}
    \footnotesize
    \begin{columns}[T,onlytextwidth]
        \begin{column}{0.56\textwidth}
            \slidejudul{Delapan bagian}
            \begin{enumerate}
                \item Dari karakter ke bit: representasi dan XOR
                \item Cipher alir dan keystream generator
                \item FSR dan LFSR
                \item Contoh terhitung LFSR
                \item RC4
                \item A5/1
                \item Trivium
                \item Salsa20 dan ChaCha20
            \end{enumerate}
        \end{column}
        \begin{column}{0.40\textwidth}
            \begin{block}{Referensi}
                \footnotesize
                Bab 5 buku ajar; Munir, \textit{IF4020 Kriptografi}, ITB,
                2026 (modul 07).
            \end{block}
        \end{column}
    \end{columns}
    \vspace{0.15cm}
    \begin{exampleblock}{Benang merah}
        \footnotesize
        Cipher alir mengganti kunci pendek yang berulang dengan aliran
        kunci acak semu yang panjang. Keamanannya bergantung pada
        generator dan umpan yang dipakai.
    \end{exampleblock}
\end{frame}

% ---------------------------------------------------- capaian pembelajaran
\begin{frame}{Capaian Pembelajaran}
    \begin{block}{Sub-CPMK 1.2 --- One-time pad dan stream cipher modern}
        \footnotesize
        Mengenal one-time pad serta stream cipher modern, yaitu representasi
        bit, operasi XOR, keystream generator, FSR, RC4, A5, Trivium,
        Salsa20, dan ChaCha20.
    \end{block}
    \vspace{0.1cm}
    {\scriptsize\itshape Pertemuan ini menuntaskan bagian stream cipher
    modern; one-time pad telah dibahas pada Minggu 6.}
    \vspace{0.15cm}
    \begin{exampleblock}{Setelah pertemuan ini, mahasiswa mampu}
        \footnotesize
        \begin{itemize}
            \item menjelaskan representasi bit, byte, dan heksadesimal;
            \item menjelaskan operasi XOR dan sifat $a\oplus a=0$;
            \item menjelaskan peran keystream generator dan umpan $U$;
            \item menjelaskan mekanisme FSR/LFSR beserta contoh keystream-nya;
            \item menjelaskan KSA/PRGA RC4, kaidah mayoritas A5/1,
                  NLFSR Trivium, dan ARX Salsa20/ChaCha20;
            \item membandingkan keempat cipher alir tersebut.
        \end{itemize}
    \end{exampleblock}
\end{frame}

% =====================================================================
\section{Dari Karakter ke Bit}
% =====================================================================

\begin{frame}{Mengapa Cipher Modern Bekerja pada Bit}
    \footnotesize
    \begin{columns}[T,onlytextwidth]
        \begin{column}{0.48\textwidth}
            \begin{block}{Cipher klasik}
                \begin{itemize}
                    \item satu huruf dipetakan ke satu huruf;
                    \item kunci berupa rangkaian huruf;
                    \item pola bahasa masih terlihat pada cipherteks.
                \end{itemize}
            \end{block}
        \end{column}
        \begin{column}{0.48\textwidth}
            \begin{block}{Cipher modern}
                \begin{itemize}
                    \item plainteks, kunci, dan cipherteks berupa rangkaian bit;
                    \item panjang kunci dinyatakan dalam bit;
                    \item operasi bekerja bit demi bit atau byte demi byte.
                \end{itemize}
            \end{block}
        \end{column}
    \end{columns}
    \vspace{0.15cm}
    \begin{block}{Yang tetap dan yang berubah}
        Substitusi dan transposisi tetap menjadi kerangka. Bedanya, operasi
        kini memakai XOR, rotasi, kompresi, ekspansi, dan penjumlahan modulo
        sehingga hubungan kunci--cipherteks lebih sulit dilacak.
    \end{block}
\end{frame}

\begin{frame}{Bit, Byte, dan Heksadesimal}
    \footnotesize
    \begin{columns}[T,onlytextwidth]
        \begin{column}{0.46\textwidth}
            \begin{block}{Satuan dasar}
                \begin{itemize}
                    \item $1$ byte $=8$ bit;
                    \item $1$ digit heksadesimal $=4$ bit;
                    \item satu byte ditulis sebagai dua digit heksadesimal;
                    \item biner dipakai saat operasi diperlihatkan langkah
                          demi langkah.
                \end{itemize}
            \end{block}
        \end{column}
        \begin{column}{0.50\textwidth}
            \begin{block}{Contoh konversi}
                \begin{itemize}
                    \item $1001\ 1101\ 0110\ 1010=\code{9D6A}$;
                    \item \code{Halo} $=$
                          \code{01001000 01100001 01101100 01101111};
                    \item \code{Indonesia emas 2045} menempati
                          $19\times 8=152$ bit.
                \end{itemize}
            \end{block}
        \end{column}
    \end{columns}
    \vspace{0.15cm}
    \begin{exampleblock}{Mengapa dipadatkan}
        Rangkaian delapan bit sulit dibaca. Heksadesimal memadatkan empat
        bit menjadi satu digit tanpa mengubah nilainya.
    \end{exampleblock}
\end{frame}

\begin{frame}{XOR: Aturan dan Tiga Sifat}
    \footnotesize
    \begin{columns}[T,onlytextwidth]
        \begin{column}{0.34\textwidth}
            \begin{block}{Tabel kebenaran}
                \centering
                \begin{tabular}{@{}ccc@{}}
                    \toprule
                    $a$ & $b$ & $a\oplus b$ \\
                    \midrule
                    0 & 0 & 0 \\
                    0 & 1 & 1 \\
                    1 & 0 & 1 \\
                    1 & 1 & 0 \\
                    \bottomrule
                \end{tabular}
            \end{block}
        \end{column}
        \begin{column}{0.62\textwidth}
            \begin{block}{Tiga sifat yang dipakai berulang}
                \begin{enumerate}
                    \item $a\oplus a=0$;
                    \item $a\oplus b=b\oplus a$ (komutatif);
                    \item $a\oplus(b\oplus c)=(a\oplus b)\oplus c$
                          (asosiatif).
                \end{enumerate}
            \end{block}
        \end{column}
    \end{columns}
    \vspace{0.15cm}
    \begin{exampleblock}{Membaca tabel}
        Hasil XOR bernilai $1$ hanya bila kedua operand berbeda. Contoh:
        $1\oplus1=0$, $1\oplus0=1$, dan $1\oplus(0\oplus1)=0$.
    \end{exampleblock}
\end{frame}

\begin{frame}{XOR Bitwise dan Sifat Membalik Diri}
    \footnotesize
    \begin{columns}[T,onlytextwidth]
        \begin{column}{0.54\textwidth}
            \begin{itemize}
                \item XOR diterapkan pada posisi bit yang bersesuaian.
                \item Contoh: $10011\oplus11001=01010$.
                \item Enkripsi: $c=p\oplus k$.
                \item Dekripsi:
                      $p'=c\oplus k=(p\oplus k)\oplus k=p$.
                \item Contoh byte: $\code{65}_{16}\oplus\code{35}_{16}
                      =\code{50}_{16}$ (karakter \code{P}); XOR sekali lagi
                      mengembalikan \code{65} (karakter \code{e}).
            \end{itemize}
            \begin{block}{Akibatnya}
                Enkripsi dan dekripsi memakai operasi yang sama. Tidak ada
                rangkaian invers terpisah.
            \end{block}
        \end{column}
        \begin{column}{0.42\textwidth}
            \centering
            \bookfig[0.98\linewidth]{05}{xor-bitwise.png}
            \figsumber{Munir, \textit{IF4020 Kriptografi}, ITB.}
        \end{column}
    \end{columns}
\end{frame}

\begin{frame}{Cipher XOR Sederhana dan Kelemahannya}
    \footnotesize
    Bila kunci lebih pendek daripada pesan, bit kunci diulang secara
    periodik:
    \vspace{0.1cm}
    \begin{center}
        \footnotesize
        \begin{tabular}{@{}l@{\quad}l@{}}
            Plainteks  & \texttt{10010010101110101010001110001} \\
            Kunci      & \texttt{11011011011011011011011011011} \\
            Cipherteks & \texttt{01001001110101110001010101010} \\
        \end{tabular}
    \end{center}
    \vspace{0.1cm}
    \begin{itemize}
        \item Kunci pada contoh hanya pola tiga bit \code{110}.
        \item Bit plainteks pada posisi $i$, $i+3$, $i+6$ memakai bit kunci
              yang sama.
        \item Pola pengulangan dapat ditemukan dengan cara yang serupa
              dengan metode Kasiski.
    \end{itemize}
    \begin{alertblock}{Kesimpulan}
        Cipher XOR sederhana \kunci{tidak aman}. Diperlukan keystream yang
        berubah sepanjang pesan, bukan kunci pendek yang berulang.
    \end{alertblock}
\end{frame}

% =====================================================================
\section{Cipher Alir dan Keystream Generator}
% =====================================================================

\begin{frame}{Prinsip Cipher Alir}
    \footnotesize
    \begin{columns}[T,onlytextwidth]
        \begin{column}{0.54\textwidth}
            \begin{itemize}
                \item Cipher alir memproses $1$ bit, $1$ byte, atau $n$ bit
                      setiap langkah.
                \item Setiap bit plainteks di-XOR dengan bit keystream pada
                      posisi yang sama:
                      \[
                          c_i=p_i\oplus k_i .
                      \]
                \item Dekripsi memakai aturan yang sama:
                      \[
                          p_i=c_i\oplus k_i .
                      \]
                \item Cipher alir langsung bekerja pada aliran data tanpa
                      menunggu data terkumpul.
            \end{itemize}
        \end{column}
        \begin{column}{0.42\textwidth}
            \centering
            \bookfig[0.98\linewidth]{05}{diagram-blok-stream-cipher.png}
            \figsumber{Munir, \textit{IF4020 Kriptografi}, ITB.}
        \end{column}
    \end{columns}
    \vspace{0.1cm}
    \begin{block}{Syarat}
        Pengirim dan penerima harus membangkitkan keystream yang identik.
    \end{block}
\end{frame}

\begin{frame}{Keystream Generator dan Umpan $U$}
    \footnotesize
    \begin{itemize}
        \item Keystream generator adalah prosedur yang sama di sisi pengirim
              dan penerima.
        \item Masukannya adalah umpan $U$; keluarannya merupakan fungsi dari
              $U$.
        \item Umpan harus sama di kedua pihak dan dijaga kerahasiaannya,
              karena kedudukannya setara dengan kunci eksternal.
        \item Keystream bersifat acak semu: dapat dibangkitkan ulang secara
              identik dan memiliki periode.
        \item Generator dapat mengeluarkan keystream dalam satuan bit, byte,
              atau blok bit.
    \end{itemize}
    \begin{exampleblock}{Konsekuensi}
        Umpan berbeda menghasilkan keystream berbeda. Bila umpan bocor,
        penyerang dapat mereproduksi seluruh keystream.
    \end{exampleblock}
\end{frame}

\begin{frame}{Galat Satu Bit Tidak Menyebar}
    \footnotesize
    \begin{itemize}
        \item Setiap bit plainteks ditentukan hanya oleh satu bit cipherteks
              dan satu bit keystream.
        \item Satu bit cipherteks yang rusak menghasilkan satu bit plainteks
              yang rusak.
        \item Tidak ada efek longsoran ke bit-bit berikutnya.
        \item Sifat ini cocok untuk aliran data kontinu dan pemrosesan
              \textit{real-time}, misalnya suara pada jaringan GSM.
    \end{itemize}
    \vspace{0.1cm}
    \begin{block}{Galat terlokalisasi}
        Galat tidak menyebar ke bit berikutnya; hanya posisi yang terkena
        yang berubah.
    \end{block}
    \begin{exampleblock}{Praktisnya}
        Transmisi suara lebih toleran terhadap satu bit yang salah daripada
        terhadap bagian data yang harus diulang seluruhnya.
    \end{exampleblock}
\end{frame}

\begin{frame}{Peta Cipher Alir yang Dibahas}
    \footnotesize
    \begin{center}
        \scriptsize
        \begin{tabularx}{\linewidth}{@{}>{\raggedright\arraybackslash}p{1.5cm}
                >{\centering\arraybackslash}p{1.0cm}
                >{\raggedright\arraybackslash}X
                >{\raggedright\arraybackslash}p{3.9cm}@{}}
            \toprule
            Cipher & Tahun & Struktur utama & Catatan \\
            \midrule
            RC4 & 1987 & larik $S$ 256 byte; KSA dan PRGA &
                pernah dipakai di SSL/TLS dan WEP/WPA \\
            \addlinespace
            A5/1 & 1989 & tiga LFSR 19, 22, dan 23 bit &
                cipher suara GSM; sudah dipecahkan \\
            \addlinespace
            Trivium & 2004 & tiga NLFSR 93, 84, dan 111 bit &
                peserta eSTREAM; ringan untuk perangkat keras \\
            \addlinespace
            Salsa20 & 2004 & matriks $4\times4$ kata 32 bit; ARX &
                peserta eSTREAM; 512 bit per panggilan \\
            \addlinespace
            ChaCha20 & 2008 & modifikasi Salsa20; ARX &
                putaran kolom dan diagonal \\
            \bottomrule
        \end{tabularx}
    \end{center}
    \vspace{0.1cm}
    \begin{block}{Rute belajar}
        RC4 dan A5/1 mewakili desain lama; Trivium dan keluarga
        Salsa20/ChaCha20 mewakili desain setelah eSTREAM.
    \end{block}
\end{frame}

% =====================================================================
\section{FSR dan LFSR}
% =====================================================================

\begin{frame}{FSR: Register Geser dengan Umpan Balik}
    \footnotesize
    \begin{columns}[T,onlytextwidth]
        \begin{column}{0.54\textwidth}
            \begin{itemize}
                \item \textit{Feedback Shift Register} (FSR) terdiri atas
                      register geser $n$ bit dan fungsi umpan balik.
                \item Setiap detak menggeser isi register satu posisi.
                \item Bit yang terlempar keluar menjadi bit keystream.
                \item Sebagian bit register diumpanbalikkan untuk mengisi
                      posisi yang kosong.
                \item Bila fungsi umpan baliknya hanya XOR, FSR itu disebut
                      LFSR.
            \end{itemize}
        \end{column}
        \begin{column}{0.42\textwidth}
            \centering
            \bookfig[0.98\linewidth]{05}{fsr-blok.png}
            \figsumber{Munir, \textit{IF4020 Kriptografi}, ITB.}
        \end{column}
    \end{columns}
    \vspace{0.1cm}
    \begin{block}{Keadaan internal}
        Isi register adalah keadaan internal generator. Umpan balik
        menentukan bit berikutnya dari keadaan itu.
    \end{block}
\end{frame}

\begin{frame}{LFSR: Umpan Balik Linier}
    \footnotesize
    \begin{columns}[T,onlytextwidth]
        \begin{column}{0.52\textwidth}
            \begin{itemize}
                \item LFSR memakai XOR saja sebagai fungsi umpan balik.
                \item Bit register ditulis $b_n b_{n-1}\dots b_1$.
                \item Setiap detak: $b_n$ pindah ke $b_{n-1}$, dan seterusnya;
                      $b_1$ keluar menjadi luaran.
                \item Posisi $b_n$ yang kosong diisi XOR bit-bit \textit{tap}.
                \item Bila semua bit dipakai:
                      $f=b_n\oplus b_{n-1}\oplus\dots\oplus b_1$.
                \item LFSR $n$ bit dapat mencapai periode $2^n-1$; keadaan
                      nol tidak dipakai.
            \end{itemize}
        \end{column}
        \begin{column}{0.44\textwidth}
            \centering
            \bookfig[0.98\linewidth]{05}{lfsr.png}
            \figsumber{Munir, \textit{IF4020 Kriptografi}, ITB.}
        \end{column}
    \end{columns}
\end{frame}

\begin{frame}{Contoh LFSR 4 Bit: Langkah demi Langkah}
    \footnotesize
    \begin{block}{Data contoh}
        Isi awal $U=\code{1111}$; umpan balik $f=b_4\oplus b_1$; register
        bergeser ke kanan. Bit luaran adalah $b_1$ sebelum register
        bergeser, dan $f$ masuk ke posisi $b_4$.
    \end{block}
    \vspace{0.1cm}
    \begin{tabular}{@{}l@{\quad}l@{\quad}l@{}}
        \toprule
        Detak & Operasi & Hasil \\
        \midrule
        $0\to1$ & $f=1\oplus1=0$; keluaran $b_1=1$ & $\code{1111}\to\code{0111}$ \\
        $1\to2$ & $f=0\oplus1=1$; keluaran $b_1=1$ & $\code{0111}\to\code{1011}$ \\
        $2\to3$ & $f=1\oplus1=0$; keluaran $b_1=1$ & $\code{1011}\to\code{0101}$ \\
        $3\to4$ & $f=0\oplus1=1$; keluaran $b_1=1$ & $\code{0101}\to\code{1010}$ \\
        \bottomrule
    \end{tabular}
    \vspace{0.15cm}
    \begin{exampleblock}{Membaca hasil}
        Setiap baris memuat operasi umpan balik, bit yang keluar, dan isi
        register baru. Empat bit luaran pertama adalah \code{1111}.
    \end{exampleblock}
\end{frame}

\begin{frame}{Tabel Lengkap 15 Detak}
    \footnotesize
    \begin{columns}[T,onlytextwidth]
        \begin{column}{0.44\textwidth}
            \centering
            \scriptsize
            \begin{tabular}{@{}cccccc@{}}
                \toprule
                $i$ & $b_4$ & $b_3$ & $b_2$ & $b_1$ & Keluaran \\
                \midrule
                0  & 1 & 1 & 1 & 1 & 1 \\
                1  & 0 & 1 & 1 & 1 & 1 \\
                2  & 1 & 0 & 1 & 1 & 1 \\
                3  & 0 & 1 & 0 & 1 & 1 \\
                4  & 1 & 0 & 1 & 0 & 0 \\
                5  & 1 & 1 & 0 & 1 & 1 \\
                6  & 0 & 1 & 1 & 0 & 0 \\
                7  & 0 & 0 & 1 & 1 & 1 \\
                8  & 1 & 0 & 0 & 1 & 1 \\
                9  & 0 & 1 & 0 & 0 & 0 \\
                10 & 0 & 0 & 1 & 0 & 0 \\
                11 & 0 & 0 & 0 & 1 & 1 \\
                12 & 1 & 0 & 0 & 0 & 0 \\
                13 & 1 & 1 & 0 & 0 & 0 \\
                14 & 1 & 1 & 1 & 0 & 0 \\
                \bottomrule
            \end{tabular}
        \end{column}
        \begin{column}{0.52\textwidth}
            \begin{itemize}
                \item Barisan luaran:
                      \code{111101011001000}.
                \item Setelah 15 detak, register kembali ke \code{1111}.
                \item Periode $=15=2^4-1$, bukan $16$.
                \item Keadaan \code{0000} tidak dipakai karena akan bernilai
                      nol selamanya.
            \end{itemize}
            \begin{block}{Mengapa periode penting}
                Keystream berulang setelah satu periode. Pesan yang lebih
                panjang dari periode memakai ulang bit awal keystream.
            \end{block}
        \end{column}
    \end{columns}
\end{frame}

\begin{frame}{Contoh Terhitung: Keystream 24 Bit}
    \footnotesize
    \begin{block}{Langkah 1 --- menyusun keystream}
        Pesan 24 bit lebih panjang daripada periode 15 bit. Keystream
        memakai 15 bit pertama, lalu 9 bit awalnya lagi.
    \end{block}
    \vspace{0.1cm}
    \begin{center}
        \scriptsize
        \begin{tabular}{@{}l@{\quad}l@{}}
            Periode 15 bit & \texttt{1111 0101 1001 000} \\
            9 bit awal     & \texttt{1111 01011} \\
            Keystream 24 bit & \texttt{1111 0101 1001 0001 1110 1011} \\
        \end{tabular}
    \end{center}
    \vspace{0.1cm}
    \[
        K=\code{111101011001000}\;\Vert\;\code{111101011}
         =\code{F591EB}.
    \]
    \begin{exampleblock}{Langkah 2 --- hasil}
        Keystream yang dipakai adalah $K=\code{F591EB}$. Sembilan bit
        terakhirnya sama dengan sembilan bit pertama karena periode
        keystream hanya 15 bit.
    \end{exampleblock}
\end{frame}

\begin{frame}{Contoh Terhitung: Enkripsi dan Dekripsi}
    \footnotesize
    \begin{block}{Data}
        Plainteks \code{Aku} $=\code{416B75}
        =\code{0100 0001 0110 1011 0111 0101}$; keystream
        $K=\code{F591EB}=\code{1111 0101 1001 0001 1110 1011}$.
    \end{block}
    \vspace{0.05cm}
    \begin{columns}[T,onlytextwidth]
        \begin{column}{0.50\textwidth}
            \begin{block}{Enkripsi per byte}
                \begin{tabular}{@{}l@{}}
                    $\code{41}\oplus\code{F5}=\code{B4}$ \\
                    $\code{6B}\oplus\code{91}=\code{FA}$ \\
                    $\code{75}\oplus\code{EB}=\code{9E}$ \\
                \end{tabular}
                \vspace{0.1cm}

                Jadi $C=\code{B4FA9E}$.
            \end{block}
        \end{column}
        \begin{column}{0.48\textwidth}
            \begin{block}{Dekripsi}
                \begin{tabular}{@{}l@{}}
                    $\code{B4}\oplus\code{F5}=\code{41}$ \\
                    $\code{FA}\oplus\code{91}=\code{6B}$ \\
                    $\code{9E}\oplus\code{EB}=\code{75}$ \\
                \end{tabular}
                \vspace{0.1cm}

                Hasilnya \code{416B75}, yaitu \code{Aku}.
            \end{block}
        \end{column}
    \end{columns}
    \vspace{0.1cm}
    \begin{exampleblock}{Pemeriksaan bit}
        Setiap posisi byte memakai operasi yang sama. Tidak ada kunci
        dekripsi khusus: keystream yang sama di-XOR-kan kembali.
    \end{exampleblock}
\end{frame}

% =====================================================================
\section{RC4}
% =====================================================================

\begin{frame}{RC4: Gambaran Umum}
    \footnotesize
    \begin{columns}[T,onlytextwidth]
        \begin{column}{0.54\textwidth}
            \begin{itemize}
                \item Dibuat Ronald Rivest pada 1987 di Laboratorium RSA.
                \item \code{RC} ditafsirkan sebagai \textit{Ron's Code} atau
                      \textit{Rivest Cipher}.
                \item Kunci paling panjang 256 byte; bila lebih pendek,
                      byte kunci diulang periodik.
                \item Keadaan internal: larik $S_0,\dots,S_{255}$ dan dua
                      pencacah $i$, $j$.
                \item Dua subproses: KSA mengacak $S$; PRGA membangkitkan
                      byte keystream.
            \end{itemize}
        \end{column}
        \begin{column}{0.42\textwidth}
            \centering
            \bookfig[0.98\linewidth]{05}{alur-rc4.png}
            \figsumber{Munir, \textit{IF4020 Kriptografi}, ITB.}
        \end{column}
    \end{columns}
    \vspace{0.1cm}
    \begin{block}{Keluaran}
        RC4 menghasilkan keystream byte demi byte tanpa batas panjang;
        setiap byte di-XOR-kan dengan satu byte plainteks.
    \end{block}
\end{frame}

\begin{frame}{KSA: Mengacak Larik $S$}
    \footnotesize
    \begin{block}{Key-Scheduling Algorithm}
        \begin{enumerate}
            \item Isi larik: $S_i=i$ untuk $i=0,1,\dots,255$.
            \item Setel $j=0$.
            \item Untuk $i=0$ sampai $255$:
                  \[
                      j\leftarrow(j+S_i+K_{i\bmod |K|})\bmod 256,
                  \]
                  lalu tukar $S_i$ dan $S_j$.
        \end{enumerate}
    \end{block}
    \vspace{0.1cm}
    \begin{itemize}
        \item Indeks kunci $i\bmod |K|$ membuat byte kunci berulang bila
              kunci lebih pendek daripada 256 byte.
        \item Hasil KSA adalah permutasi larik $S$ yang ditentukan oleh
              kunci rahasia.
    \end{itemize}
    \begin{exampleblock}{Peran KSA}
        KSA menyiapkan keadaan awal. Tanpa KSA, PRGA akan memakai larik
        berurutan yang sama untuk setiap kunci.
    \end{exampleblock}
\end{frame}

\begin{frame}{PRGA: Membangkitkan Byte Keystream}
    \footnotesize
    \begin{columns}[T,onlytextwidth]
        \begin{column}{0.54\textwidth}
            \begin{block}{Pseudo-random generation algorithm}
                \begin{enumerate}
                    \item Setel $i=j=0$.
                    \item Untuk setiap byte:
                          \begin{align*}
                              i&\leftarrow(i+1)\bmod 256,\\
                              j&\leftarrow(j+S_i)\bmod 256,
                          \end{align*}
                          tukar $S_i$ dan $S_j$, lalu
                          \[
                              t\leftarrow(S_i+S_j)\bmod 256,\qquad
                              u\leftarrow S_t .
                          \]
                    \item Byte keystream $u$ di-XOR-kan dengan plainteks:
                          $C=P\oplus u$.
                \end{enumerate}
            \end{block}
        \end{column}
        \begin{column}{0.42\textwidth}
            \centering
            \bookfig[0.98\linewidth]{05}{rc4-prga.png}
            \figsumber{Munir, \textit{IF4020 Kriptografi}, ITB.}
        \end{column}
    \end{columns}
    \vspace{0.05cm}
    \begin{block}{Dekripsi}
        Prosedur PRGA yang sama dipakai dengan kunci yang sama.
    \end{block}
\end{frame}

\begin{frame}{Keamanan RC4}
    \footnotesize
    \begin{itemize}
        \item Byte keystream awal berkorelasi tinggi dengan beberapa byte
              awal kunci.
        \item Implementasi yang hati-hati membuang 256 sampai 512 byte
              keystream pertama.
        \item RC4 lemah terhadap serangan \textit{flip-bit} dan serangan
              \textit{stream} lain.
        \item Keamanannya dapat dipecahkan dalam hitungan jam atau hari.
        \item RFC 7465 (Februari 2015) melarang pemakaian RC4 di TLS.
        \item Varian perbaikan: Spritz, RC4A, VMPC, dan RC4+.
    \end{itemize}
    \begin{alertblock}{Pelajaran}
        Panjang kunci bukan jaminan keamanan. RC4 menunjukkan bahwa cipher
        cepat tetap dapat kehilangan keamanannya.
    \end{alertblock}
\end{frame}

% =====================================================================
\section{A5/1}
% =====================================================================

\begin{frame}{A5/1: Cipher Alir untuk GSM}
    \footnotesize
    \begin{columns}[T,onlytextwidth]
        \begin{column}{0.54\textwidth}
            \begin{itemize}
                \item A5 adalah cipher alir untuk transmisi suara pada
                      jaringan GSM.
                \item Satu frame GSM panjangnya 228 bit dan dikirim setiap
                      4,6 milidetik.
                \item Kunci eksternalnya adalah kunci sesi 64 bit.
                \item A5/1 adalah versi kuat yang dipakai di Eropa; A5/2
                      adalah versi ekspor yang lebih lemah.
                \item Algoritma awalnya rahasia, lalu terbongkar melalui
                      rekayasa balik pada 1994.
            \end{itemize}
        \end{column}
        \begin{column}{0.42\textwidth}
            \centering
            \bookfig[0.98\linewidth]{05}{a51-lfsr.png}
            \figsumber{Munir, \textit{IF4020 Kriptografi}, ITB.}
        \end{column}
    \end{columns}
    \vspace{0.1cm}
    \begin{block}{Keluaran per frame}
        Setiap frame memakai keystream 228 bit yang di-XOR-kan dengan
        bit-bit frame pesan.
    \end{block}
\end{frame}

\begin{frame}{Tiga LFSR dan Kaidah Mayoritas}
    \footnotesize
    \begin{columns}[T,onlytextwidth]
        \begin{column}{0.50\textwidth}
            \centering
            \scriptsize
            \begin{tabular}{@{}lccc@{}}
                \toprule
                Register & Panjang & Kendali & Umpan balik \\
                \midrule
                $R_1$ & 19 & 8  & 13, 16, 17, 18 \\
                $R_2$ & 22 & 10 & 20, 21 \\
                $R_3$ & 23 & 10 & 7, 20, 21, 22 \\
                \bottomrule
            \end{tabular}
            \vspace{0.1cm}

            Total $19+22+23=64$ bit.
        \end{column}
        \begin{column}{0.46\textwidth}
            \begin{itemize}
                \item Bit kendali tiap register dibandingkan dengan bit
                      mayoritas dari ketiganya.
                \item Register yang bit kendalinya sama dengan mayoritas
                      akan didetak.
                \item Peluang sebuah register didetak adalah $3/4$; biasanya
                      dua atau tiga register bergeser.
                \item Luaran:
                      \[
                          z_t=R_1[18]\oplus R_2[21]\oplus R_3[22].
                      \]
            \end{itemize}
        \end{column}
    \end{columns}
\end{frame}

\begin{frame}{Inisialisasi dan Pembangkitan Keystream}
    \footnotesize
    \begin{enumerate}
        \item \textbf{Muat kunci.} Ketiga register diisi nol, lalu 64 putaran.
              Pada putaran ke-$i$, bit $K[i]$ di-XOR-kan ke LSB setiap
              register sebelum register didetak.
        \item \textbf{Muat nomor frame.} Lakukan 22 putaran tambahan; bit
              $F_n[i]$ di-XOR-kan ke LSB setiap register.
        \item \textbf{Pemanasan.} Detak 100 putaran dengan kaidah mayoritas;
              bit luarannya dibuang.
        \item \textbf{Bangkitkan keystream.} Detak 228 putaran dengan kaidah
              mayoritas; setiap putaran menghasilkan satu bit menurut
              $z_t=R_1[18]\oplus R_2[21]\oplus R_3[22]$.
    \end{enumerate}
    \vspace{0.1cm}
    \begin{block}{Urutan kerja}
        Tahap 1 dan 2 menyiapkan keadaan, tahap 3 mencampur kunci dengan
        nomor frame, dan tahap 4 menghasilkan 228 bit keystream.
    \end{block}
\end{frame}

\begin{frame}{Keamanan A5/1}
    \footnotesize
    \begin{itemize}
        \item Tujuan kriptanalisis adalah menemukan kunci sesi dari
              keystream awal yang diketahui.
        \item Serangan tergolong \textit{known-plaintext attack}.
        \item Ross Anderson (1994) menerka 42 bit isi $R_1$ dan $R_2$,
              lalu menurunkan 23 bit $R_3$.
        \item Pekerjaan itu memerlukan lebih dari sebulan pada sebuah PC.
        \item Pada tahun-tahun berikutnya, kunci dapat ditemukan hanya
              dalam beberapa menit.
    \end{itemize}
    \begin{alertblock}{Status}
        A5/1 sudah berhasil dikriptanalisis. LFSR linear dengan keadaan
        64 bit tidak cukup untuk menjaga percakapan GSM.
    \end{alertblock}
\end{frame}

% =====================================================================
\section{Trivium}
% =====================================================================

\begin{frame}{Trivium: Tiga NLFSR}
    \footnotesize
    \begin{columns}[T,onlytextwidth]
        \begin{column}{0.54\textwidth}
            \begin{itemize}
                \item Diajukan Christophe De Cannière dan Bart Preneel pada
                      kompetisi eSTREAM 2004--2008.
                \item Tiga \textit{nonlinear LFSR} (NLFSR): $A=93$ bit,
                      $B=84$ bit, dan $C=111$ bit.
                \item Total keadaan internal $93+84+111=288$ bit.
                \item Gerbang AND membuat strukturnya nonlinier, berbeda
                      dari LFSR murni.
                \item Luaran:
                      $s_i=a_i\oplus b_i\oplus c_i$.
            \end{itemize}
        \end{column}
        \begin{column}{0.42\textwidth}
            \centering
            \bookfig[0.98\linewidth]{05}{trivium-sirkuit.png}
            \figsumber{Munir, \textit{IF4020 Kriptografi}, ITB.}
        \end{column}
    \end{columns}
    \vspace{0.1cm}
    \begin{block}{Implementasi perangkat keras}
        Trivium hanya memerlukan 288 sel register, tiga gerbang AND, dan
        tujuh gerbang XOR.
    \end{block}
\end{frame}

\begin{frame}{Trivium: Titik Umpan Balik}
    \footnotesize
    \begin{center}
        \scriptsize
        \begin{tabular}{@{}lcccc@{}}
            \toprule
            Register & Panjang & Umpan balik & \textit{Feedforward} & Masukan AND \\
            \midrule
            $A$ & 93 bit  & 69 & 66 & 91, 92 \\
            $B$ & 84 bit  & 78 & 69 & 82, 83 \\
            $C$ & 111 bit & 87 & 66 & 109, 110 \\
            \bottomrule
        \end{tabular}
    \end{center}
    \vspace{0.1cm}
    \begin{itemize}
        \item Umpan balik, \textit{feedforward}, dan masukan gerbang AND
              menentukan pembaruan nonlinier setiap register.
        \item Luaran tetap XOR dari luaran ketiga register:
              $s_i=a_i\oplus b_i\oplus c_i$.
        \item Panjang dan titik tap yang berbeda membuat ketiga register
              bergerak dengan pola yang tidak sama.
    \end{itemize}
    \begin{block}{Membaca tabel}
        Kolom \textit{feedforward} dan masukan AND adalah posisi bit yang
        diambil dari register, bukan nilai bitnya.
    \end{block}
\end{frame}

\begin{frame}{Trivium: Inisialisasi dan Pemanasan}
    \footnotesize
    \begin{itemize}
        \item \textit{Initialization vector} 80 bit dimasukkan ke register
              $A$.
        \item Kunci 80 bit dimasukkan ke register $B$.
        \item Register $C$ diisi nol, kecuali tiga bit terakhir:
              $c_{109}=c_{110}=c_{111}=1$.
        \item Ketiga NLFSR didetak $4\times288=1152$ kali tanpa
              menghasilkan luaran.
        \item Pemanasan mencampur kunci dan IV sehingga pengaruhnya
              tersebar ke seluruh register.
    \end{itemize}
    \begin{exampleblock}{Mengapa ringan}
        Keadaan 288 bit dan gerbang sederhana membuat Trivium cocok untuk
        \textit{lightweight cryptography} pada perangkat dengan sumber
        daya terbatas.
    \end{exampleblock}
\end{frame}

% =====================================================================
\section{Salsa20 dan ChaCha20}
% =====================================================================

\begin{frame}{Salsa20: Keadaan $4\times4$}
    \footnotesize
    \begin{columns}[T,onlytextwidth]
        \begin{column}{0.54\textwidth}
            \begin{itemize}
                \item Dikembangkan Daniel J. Bernstein pada 2005 dan
                      diajukan pada eSTREAM 2004--2008.
                \item Kunci rahasia 256 bit; \textit{nonce} 64 bit; pencacah
                      64 bit yang dimulai dari nol.
                \item Satu pemanggilan menghasilkan keystream 512 bit
                      atau 64 byte.
                \item Keadaan internal: matriks $4\times4$ berisi 16 kata.
                \item Satu kata 32 bit, sehingga ukurannya
                      $16\times32=512$ bit.
                \item \textit{Nonce} dan pencacah tidak perlu dirahasiakan;
                      kunci harus tetap rahasia.
            \end{itemize}
        \end{column}
        \begin{column}{0.42\textwidth}
            \centering
            \bookfig[0.98\linewidth]{05}{salsa20-state.png}
            \figsumber{Munir, \textit{IF4020 Kriptografi}, ITB.}
        \end{column}
    \end{columns}
\end{frame}

\begin{frame}{Salsa20: Putaran dan \textit{Quarter Round}}
    \footnotesize
    \begin{columns}[T,onlytextwidth]
        \begin{column}{0.50\textwidth}
            \begin{block}{\textit{Quarter round} (QR)}
                \begin{align*}
                    b&\leftarrow b\oplus(a+d)\lll 7,\\
                    c&\leftarrow c\oplus(b+a)\lll 9,\\
                    d&\leftarrow d\oplus(c+b)\lll 13,\\
                    a&\leftarrow a\oplus(d+c)\lll 18.
                \end{align*}
            \end{block}
            {\scriptsize $+$ adalah penjumlahan modulo $2^{32}$;
            $\lll k$ adalah rotasi kiri $k$ bit.}
        \end{column}
        \begin{column}{0.46\textwidth}
            \begin{itemize}
                \item Satu putaran ganda terdiri atas putaran kolom dan
                      putaran baris.
                \item Salsa20 melakukan 10 putaran ganda, yaitu 20 putaran.
                \item Setelah itu hasilnya dijumlahkan dengan keadaan awal:
                      \[
                          y_j=x'_j+x_j\bmod 2^{32}.
                      \]
                \item Pencacah bertambah satu untuk setiap blok 512 bit;
                      \textit{nonce} tetap.
            \end{itemize}
        \end{column}
    \end{columns}
    \vspace{0.1cm}
    \begin{block}{Difusi}
        Perubahan kecil pada kunci, \textit{nonce}, atau pencacah
        menghasilkan perubahan besar pada keystream.
    \end{block}
\end{frame}

\begin{frame}{ChaCha20: Tiga Perbedaan dari Salsa20}
    \footnotesize
    \begin{columns}[T,onlytextwidth]
        \begin{column}{0.48\textwidth}
            \begin{enumerate}
                \item Tata letak keadaan awal berbeda.
                \item Bentuk QR berbeda, dengan rotasi 16, 12, 8, dan 7 bit.
                \item ChaCha memakai putaran kolom dan putaran diagonal.
            \end{enumerate}
            \vspace{0.1cm}
            \begin{block}{Keluarga ChaCha}
                ChaCha20 adalah anggota keluarga ChaCha bersama ChaCha6 dan
                ChaCha7. Tujuannya adalah difusi yang lebih baik dengan
                kecepatan dan kesederhanaan ARX.
            \end{block}
        \end{column}
        \begin{column}{0.48\textwidth}
            \begin{block}{QR ChaCha20}
                \scriptsize
                \begin{align*}
                    a&\leftarrow a+b, &
                    d&\leftarrow d\oplus a, &
                    d&\leftarrow d\lll 16,\\
                    c&\leftarrow c+d, &
                    b&\leftarrow b\oplus c, &
                    b&\leftarrow b\lll 12,\\
                    a&\leftarrow a+b, &
                    d&\leftarrow d\oplus a, &
                    d&\leftarrow d\lll 8,\\
                    c&\leftarrow c+d, &
                    b&\leftarrow b\oplus c, &
                    b&\leftarrow b\lll 7.
                \end{align*}
            \end{block}
        \end{column}
    \end{columns}
    \vspace{0.1cm}
    \begin{exampleblock}{ARX}
        Seluruh operasi hanya penjumlahan (\textit{add}), rotasi
        (\textit{rotate}), dan XOR. Tidak ada perkalian atau tabel
        substitusi.
    \end{exampleblock}
\end{frame}

\begin{frame}{ChaCha20: Putaran dan Keunggulan}
    \footnotesize
    \begin{columns}[T,onlytextwidth]
        \begin{column}{0.50\textwidth}
            \begin{block}{Pemilihan indeks kata}
                \scriptsize
                \begin{tabular}{@{}ll@{}}
                    \toprule
                    Putaran kolom & Putaran diagonal \\
                    \midrule
                    QR(0, 4, 8, 12)  & QR(0, 5, 10, 15) \\
                    QR(1, 5, 9, 13)  & QR(1, 6, 11, 12) \\
                    QR(2, 6, 10, 14) & QR(2, 7, 8, 13) \\
                    QR(3, 7, 11, 15) & QR(3, 4, 9, 14) \\
                    \bottomrule
                \end{tabular}
            \end{block}
            {\scriptsize Empat QR per putaran kolom dan empat QR per putaran
            diagonal.}
        \end{column}
        \begin{column}{0.46\textwidth}
            \begin{itemize}
                \item ARX cepat di perangkat lunak dan mudah diimplementasikan.
                \item Setiap blok keystream bergantung pada kunci,
                      \textit{nonce}, dan pencacah.
                \item Blok dapat dihitung secara independen dan paralel.
                \item ChaCha20 banyak dipakai dalam protokol modern.
            \end{itemize}
        \end{column}
    \end{columns}
    \vspace{0.1cm}
    \begin{block}{Inti perbandingan}
        Salsa20 dan ChaCha20 sama-sama menghasilkan blok 512 bit dengan
        kunci 256 bit. ChaCha20 mengubah tata letak keadaan dan pola QR
        agar difusi per putaran lebih baik.
    \end{block}
\end{frame}

% =====================================================================
\section{Rangkuman dan Penutup}
% =====================================================================

\begin{frame}{Rangkuman (1/2)}
    \footnotesize
    \begin{itemize}
        \item Cipher modern bekerja pada bit dan byte; XOR menjadi operasi
              dasar karena $a\oplus a=0$ membuat enkripsi dan dekripsi
              memakai fungsi yang sama.
        \item Cipher alir mengenkripsi per bit, byte, atau $n$ bit dengan
              $c_i=p_i\oplus k_i$.
        \item Keamanan cipher alir bergantung pada keystream generator:
              generator menerima umpan $U$ dan menghasilkan keystream acak
              semu yang dapat dibangkitkan ulang.
        \item FSR menggeser register dan mengumpanbalikkan sebagian bit;
              LFSR membatasi umpan balik pada XOR.
        \item LFSR $n$ bit dapat mencapai periode $2^n-1$. Contoh 4 bit
              dengan $U=\code{1111}$ menghasilkan periode 15 bit.
    \end{itemize}
\end{frame}

\begin{frame}{Rangkuman (2/2)}
    \footnotesize
    \begin{center}
        \scriptsize
        \begin{tabularx}{\linewidth}{@{}>{\raggedright\arraybackslash}p{1.4cm}
                >{\raggedright\arraybackslash}p{3.1cm}
                >{\raggedright\arraybackslash}p{1.9cm}
                >{\raggedright\arraybackslash}X@{}}
            \toprule
            Cipher & Struktur & Keluaran & Catatan \\
            \midrule
            RC4 & larik $S$ 256 byte; KSA dan PRGA & byte &
                byte awal berkorelasi dengan kunci; tidak aman \\
            \addlinespace
            A5/1 & tiga LFSR 19/22/23 bit & bit per detak &
                kaidah mayoritas; sudah dipecahkan \\
            \addlinespace
            Trivium & tiga NLFSR 93/84/111 bit & bit &
                pemanasan 1152 detak; ringan \\
            \addlinespace
            Salsa20 & $4\times4$ kata 32 bit & 512 bit per panggilan &
                20 putaran; ARX \\
            \addlinespace
            ChaCha20 & $4\times4$ kata 32 bit & 512 bit per panggilan &
                QR 16/12/8/7; kolom dan diagonal \\
            \bottomrule
        \end{tabularx}
    \end{center}
    \vspace{0.1cm}
    \begin{block}{Benang merah}
        RC4 dan A5/1 menunjukkan bahwa desain lama sudah kedaluwarsa.
        Trivium menambah nonlinearitas melalui gerbang AND, sedangkan
        Salsa20/ChaCha20 memakai ARX untuk difusi yang cepat.
    \end{block}
\end{frame}

\begin{frame}{Penutup}
    \begin{block}{Tiga hal yang harus dibawa dari pertemuan ini}
        \begin{enumerate}
            \item Cipher alir adalah XOR plainteks dengan keystream acak
                  semu; keamanannya bergantung pada generator dan umpan.
            \item LFSR adalah generator paling sederhana; periode maksimumnya
                  $2^n-1$, dan contoh 4 bit menghasilkan 15 bit keystream.
            \item RC4 dan A5/1 sudah dipecahkan; Trivium memakai NLFSR dan
                  gerbang AND, sedangkan Salsa20/ChaCha20 memakai ARX.
        \end{enumerate}
    \end{block}
    \vspace{0.3cm}
    \begin{exampleblock}{Pertemuan berikutnya --- Minggu 8}
        \textbf{Minggu 8 --- UTS}\\
        Ujian Tengah Semester: CPMK-2 serta CPMK-1 (cipher klasik, OTP,
        dan stream cipher).
    \end{exampleblock}
\end{frame}

\end{document}
